% Generated draft: requires mathematical and source review.
\documentclass[11pt]{article}
\usepackage[margin=25mm]{geometry}
\usepackage{fontspec}
\setmainfont{Latin Modern Roman}
\usepackage{amsmath,amssymb,mathtools,unicode-math}
\setmathfont{Latin Modern Math}
\usepackage{enumitem,needspace,xurl,hyperref}
\setlength{\parindent}{0pt}
\setlength{\parskip}{5pt}
\setlength{\emergencystretch}{3em}
\newcommand{\N}{\mathbb{N}}
\newcommand{\Z}{\mathbb{Z}}
\newcommand{\Q}{\mathbb{Q}}
\newcommand{\R}{\mathbb{R}}
\newcommand{\C}{\mathbb{C}}
\newcommand{\F}{\mathbb{F}}
\newcommand{\A}{\mathbb{A}}
\newcommand{\PP}{\mathbb{P}}
\newcommand{\E}{\mathbb{E}}
\newcommand{\eps}{\varepsilon}
\newcommand{\dd}{\,\mathrm{d}}
\begin{document}
\section*{UnsolvedMath ID: 700197}
\textbf{Source problem code:} catalogue entry
\subsection{Definitions and mathematical statement}
For classes \(\mathcal A\) and \(\mathcal B\) of groups, an extension-product convention is
\[
\mathcal A\mathcal B=\{G:\text{there exists }N\trianglelefteq G\text{ with }N\in\mathcal A\text{ and }G/N\in\mathcal B\}.
\]
The supplied excerpt refers to products of classes without defining the convention. Let \(\mathfrak H\) denote the collection of hereditary radical classes intended by the source. The source asserts that these classes form a semigroup under their product. It does not supply the axioms defining a radical class, the precise meaning of hereditary, or the definition of an indecomposable element.

For a prime \(p\), define
\[
\mathcal L_p=\{G:\langle X\rangle\text{ is finite for every finite }X\subseteq G,\ \text{and every element of }G\text{ has order }p^n\text{ for some }n\in\Z_{\geq0}\}.
\]
Thus \(\mathcal L_p\) is the class of locally finite \(p\)-groups.

\textbf{Classification problem attributed to B. I. Plotkin.} Determine all indecomposable elements of \(\mathfrak H\). In particular, determine those indecomposable elements \(\mathcal R\in\mathfrak H\) for which \(\mathcal R\subseteq\mathcal L_p\).

\textbf{Unresolved formulation.} One possible interpretation of indecomposability, using the displayed extension product, is
\[
\nexists\mathcal A,\mathcal B\in\mathfrak H:\quad
\mathcal R=\mathcal A\mathcal B,\qquad
\mathcal A\subsetneq\mathcal R,\qquad
\mathcal B\subsetneq\mathcal R.
\]
This is a provisional interpretation, not a definition established by the supplied record. The source must clarify whether decomposition means a product of two strictly smaller classes, a product with neither factor the identity, or another condition. The intended radical and heredity axioms must also be supplied before the classification target is fully specified.
\subsection{Short English statement}
Plotkin asks for a classification of the indecomposable hereditary radical classes of groups, especially those contained in the locally finite \(p\)-groups. The supplied excerpt leaves the radical axioms and the meaning of indecomposable unspecified, so the formal interpretation remains provisional.
\subsection{Sources}
\begin{itemize}
\item Supplied local UnsolvedMath registry record: permanent ID \(700197\), exact source code \texttt{catalogue entry}, title \emph{Kourovka Notebook Problem 3.48}; upstream revision \texttt{d371ee917ac9647fbb06cd6fc529deec6f3dda82}.
\item Supplied attribution: B. I. Plotkin, proposed in 1969; \emph{The Kourovka Notebook (Unsolved Problems in Group Theory)}, original problem 3.48, 3rd issue (1969). The record identifies the consulted edition as the 21st issue edition (2026), \texttt{arXiv:1401.0300v48}. This bibliographic attribution is supplied evidence and has not been independently checked.
\item The registry labels the problem \texttt{open}. This is a supplied status label, not an independently verified assessment of the literature. No separate literature sources or literature-verification date are supplied.
\end{itemize}
\end{document}
